% Conversão fiel de chapter-02-text-h-0003; origem RELATORIO_ATUAL.odt.
\hypertarget{_Toc45627143}{}
\hypertarget{__RefHeading___Toc4802_3790000799}{}
\section{Introdução}
\label{ages-i:text-h-0003}

\indent \originaltab{}O projeto Easy Choose consiste na criação de uma plataforma que auxilia consumidores na escolha dos melhores suplementos para suas dietas por meio de um fluxograma apresentado em um tablet disponível em lojas de varejo. A interface do tablet mostra os melhores suplementos disponíveis na loja, baseando-se nas informações inseridas pelo usuário.\par

\indent \originaltab{}Além disso, o projeto inclui uma interface para os lojistas, permitindo o cadastro de tablets e produtos em suas lojas, e uma interface para os stakeholders, que podem gerenciar os lojistas e acessar outras métricas geradas pelos tablets.\originalspace{}\par

\indent \originaltab{}O período de execução ocorreu de 28 de fevereiro de 2024 a 12 de junho de 2024, sob a orientação do professor Daniel Antônio Callegari. A equipe está retratada na Figura~\ref{ages-i:draw-image-0001}.\par

\memorialfigure{imagens/original/asset-0012.jpg}{\textbf{Figura 1 –} Foto do time – AGES I}{ages-i:draw-image-0001}{Fonte: pendente de identificação}{0.92}
